\begin{EnglishAbstract}
	
	Duchenne Muscular Dystrophy (DMD) is one of the most prevalent and severe genetic neuromuscular disorders in children, characterized by progressive skeletal muscle degeneration and gradual loss of motor function. Quantitative and reliable assessment of motor performance in patients with DMD plays a crucial role in monitoring disease progression, evaluating treatment effectiveness, and designing rehabilitation programs. One of the most widely used clinical tools for functional evaluation is the North Star Ambulatory Assessment (NSAA), a 17-item rating scale in which each task is scored from 0 to 2 based on clinical observation. Despite its clinical value, the qualitative nature of NSAA and its reliance on subjective judgment may introduce inter-rater variability and limit repeatability.
	
	The objective of this study is to design and develop an inertial sensor–based system for quantitative evaluation of motor function indices in patients with DMD, aiming to provide a data-driven complementary framework for clinical scoring. In this approach, multiple inertial measurement units (IMUs), consisting of accelerometers and gyroscopes, were attached to key body segments to record linear acceleration and angular velocity during selected NSAA activities. The raw sensor data were transmitted to the MATLAB environment, where signal preprocessing—including noise reduction, digital filtering, and temporal normalization—was performed. Subsequently, relevant time-domain and frequency-domain features were extracted to construct a structured dataset for further analysis.
	
	Machine learning algorithms were employed to estimate the score of each motor task based on the extracted features, and the predicted scores were compared with clinical evaluations. In the preliminary phase of the study, approximately 12 patients with DMD were clinically observed in collaboration with the medical team, providing valuable insight into motor behavior patterns and practical system constraints. The findings indicate that the proposed inertial sensor–based and data-driven framework can serve as a complementary tool to conventional clinical assessment by enhancing objectivity, accuracy, and repeatability in motor performance evaluation.
	
	The proposed system has the potential to support longitudinal patient monitoring, remote assessment applications, and improved clinical decision-making in future developments.
	
	\enkeywords{Duchenne Muscular Dystrophy, Inertial Sensors, Motion Analysis, Machine Learning, NSAA, Motor Function Assessment}
	
\end{EnglishAbstract}
